\begin{abstract}
Evaluating generalization in Large Language Models (LLMs) is increasingly difficult due to widespread data contamination in modern web-scale training corpora. We investigate whether models trained on sparse, historically-cutoff data provide a more scientifically rigorous testbed for measuring generalization. Specifically, we test the \hypo hypothesis: that the lack of modern concepts in a model's training data creates a clear, binary-like signal for generalization when those concepts are introduced via in-context learning. We compare \talkie, a 13B model with a strict 1930 cutoff, against a modern baseline (\talkieweb) on tasks that were non-existent in 1930, such as Python coding (\humaneval). Our results show that while \talkie fails entirely in a \zeroshot setting (0\% success), it achieves a 33\% success rate on simple logic tasks after just 3 shots, representing a ``pure'' logic leap. Furthermore, we identify a sharp \cliff in performance on subjects dominated by post-1930 knowledge. We conclude that historically-limited training data simplifies the detection of generalization by eliminating the noise of memorized knowledge, making \talkie a valuable instrument for studying the mechanics of machine reasoning.
\end{abstract}
